% vol_reb_decomposition.tex
% Derivation and interpretation of the size-vs-level decomposition of the
% variance of consumption-share growth. The decomposition is computed by
% share_var_decomposition() in reg_cons_insurance.py, called at the end of
% the experiment scripts (alimony experiment.py, WLP_experiment.py,
% tax_experiment.py).
\documentclass[11pt]{article}
\usepackage{amsmath,amssymb}
\usepackage[margin=1.1in]{geometry}
\newcommand{\E}{\mathbb{E}}
\newcommand{\ren}{\mathcal{R}}

\title{Decomposing the variance of consumption-share growth:\\ size versus level of renegotiations}
\author{}
\date{}

\begin{document}
\maketitle

\section{Object and notation}

Let $s^g_{it}$ denote spouse $g$'s share of private consumption within
couple $i$ at age $t$ (for the wife, $s^w = c^w/(c^w+c^m)$; for the husband,
$s^m = 1-s^w$), and let
\[
y^g_{it} \;\equiv\; \Delta \log s^g_{it} \;=\; \log s^g_{i,t+1}-\log s^g_{it}
\]
be its log growth over married cells. The object of interest is the
(uncentered) volatility $\E\!\left[(y^g)^2\right]$, the ``share component''
of individual private-consumption volatility in the policy-experiment
figures. In the model, $s^g$ is a monotone function of the Pareto weight
alone, so $y^g$ is (up to numerical interpolation error) nonzero only in
periods in which the weight is renegotiated. Let $\ren$ denote the set of
renegotiation cells (the bargaining weight differs from its lag) and
$f \equiv \Pr(\ren)$ the renegotiation frequency.

\section{Derivation}

Write the log change as the product of the share change in \emph{levels}
and an effective inverse share:
\[
y^g_{it} \;=\; \Delta S^g_{it}\cdot M^g_{it},
\qquad
\Delta S^g_{it} \equiv s^g_{i,t+1}-s^g_{it},
\qquad
M^g_{it} \equiv \frac{\Delta \log s^g_{it}}{\Delta S^g_{it}} .
\]
This is an identity, not an approximation. By the mean value theorem applied
to $\log(\cdot)$, $M^g_{it} = 1/\tilde{s}^g_{it}$ for some
$\tilde{s}^g_{it}$ between $s^g_{it}$ and $s^g_{i,t+1}$: $M$ is the exact
``$1/\text{share}$'' \emph{at which the change is evaluated}. Splitting
$\E[(y^g)^2]$ into renegotiation and non-renegotiation cells and factoring
the renegotiation part,
\begin{equation}
\E\!\left[(y^g)^2\right]
\;=\;
\underbrace{f \cdot
\E\!\left[\Delta S^2 \,\middle|\, \ren\right]\cdot
\E\!\left[(M^g)^2 \,\middle|\, \ren\right]\cdot
R^g}_{\text{renegotiation part}}
\;+\;
\underbrace{(1-f)\,\E\!\left[(y^g)^2 \,\middle|\, \neg\ren\right]}_{\text{numerical noise (reported, small)}},
\label{eq:dec}
\end{equation}
where
\[
R^g \;\equiv\;
\frac{\E\!\left[\Delta S^2 (M^g)^2 \mid \ren\right]}
     {\E\!\left[\Delta S^2 \mid \ren\right]\,\E\!\left[(M^g)^2 \mid \ren\right]}
\]
is the correction that turns the product of averages into the average of the
product. Equation~\eqref{eq:dec} is exact; the four factors of the
renegotiation part are log-additive, so changes across policy variants can
be attributed factor by factor.

A structural implication built into the decomposition: since
$\Delta S^m = -\Delta S^w$, the frequency $f$ and the size factor
$\E[\Delta S^2\mid\ren]$ are \emph{identical for the two spouses}. Any
difference between the wife's and the husband's share volatility---in level
or in its response to policy---must run through the level factor
$\E[(M^g)^2\mid\ren]$ and the interaction $R^g$.

\section{Interpretation of the factors}

\begin{itemize}
\item[{$f$:}] \textbf{how often} the bargaining weight moves. A policy that
  tightens or relaxes participation constraints shows up here.
\item[{$\E[\Delta S^2\mid\ren]$:}] \textbf{how large} renegotiations are in
  \emph{levels} of the share---bargaining risk with the proportional
  (log) scaling stripped out. Common to both spouses.
\item[{$\E[(M^g)^2\mid\ren]$:}] \textbf{where in the share distribution} the
  events land. A given level change is proportionally large for the spouse
  whose share is small: for a share near $0.3$ this factor is roughly
  $(1/0.3)^2 \approx 11$, against $(1/0.7)^2 \approx 2$ for the other
  spouse. Because it is computed event by event, it captures composition:
  if a policy relocates renegotiations toward couples with more equal
  shares, this factor falls for the low-share spouse even if the mean share
  moves little.
\item[{$R^g$:}] the \textbf{pairing of size and location}. $R>1$ means the
  large share changes happen precisely where the spouse's share is small
  (damage compounds); $R<1$ means they happen where the share is large
  (damage cushioned); $R=1$ means size and location are unrelated and the
  size and level factors can be discussed independently. $R$ is a
  bookkeeping term, not a separate economic channel: if it is close to one,
  it should not be interpreted.
\end{itemize}

The non-renegotiation term in~\eqref{eq:dec} collects the microscopic share
movements that occur off renegotiation cells (interpolation of the
intra-period allocation over the weight/resources grids). It is reported as
a share of $\E[(y^g)^2]$ and should be negligible; a non-negligible value
would indicate that measured share volatility partly reflects numerical
noise rather than bargaining.

\section{Implementation}

The decomposition is computed by \texttt{share\_var\_decomposition()} in
\texttt{reg\_cons\_insurance.py} and printed at the end of each policy
experiment (\texttt{alimony experiment.py}, \texttt{WLP\_experiment.py},
\texttt{tax\_experiment.py}), one row per policy variant and spouse, on the
same married cells as the corresponding volatility figures. Column mapping:
\texttt{E[y2]x100} $=\E[(y^g)^2]\times 100$; \texttt{freq} $=f$;
\texttt{E[dS2|r]x100} $=\E[\Delta S^2\mid\ren]\times 100$;
\texttt{E[M2|r]} $=\E[(M^g)^2\mid\ren]$; \texttt{R} $=R^g$;
\texttt{nonren\%} $=$ the non-renegotiation share of $\E[(y^g)^2]$ in
percent. The exactness of~\eqref{eq:dec} is verified numerically at every
call (an \texttt{identity gap} warning is printed if it fails).

\end{document}
